\section{Conclusion}
This study has successfully designed, implemented, and validated a lightweight, standards-compliant Modbus TCP threat monitoring framework by integrating rule-based deep packet inspection (DPI) via Suricata with centralized security information and event management (SIEM) via Graylog. The proposed architecture systematically aligns technical detections with the foundational requirements of ISA/IEC 62443-3-3, the tactical stages of the ICS Cyber Kill Chain, and adversary behaviors mapped in MITRE ATT\&CK for ICS.

Experimental evaluations demonstrate the practical efficiency and reliability of the developed defense pipeline. Across all tested scenarios, the system maintained an average end-to-end alerting latency of under 10 seconds. Under high-throughput loading during peak-rate stress testing, the passive pipeline successfully ingested a sustained volume of up to 29,128.00 Packets Per Second (PPS) (approximately 13.36 Mbps) while maintaining a 0.00\% packet drop rate at the Suricata kernel socket and a 0.00\% log loss rate in the SIEM layer. Furthermore, supplementary validation on third-party datasets (CICModbusDataset2023 and 4SICS Geek Lounge) confirmed the portability and detection performance of the custom rulesets within heterogeneous environments. This open-source framework provides a cost-effective solution for industrial organizations to establish continuous security visibility and satisfy compliance standards without actively interfering with real-time physical control processes.

\section{Future Work}
Future research can expand and refine this framework through the following directions:

\begin{enumerate}
    \item \textbf{Protocol and Semantic Expansion:} Extending the DPI-SIEM co-detection rules to other widely used industrial communication protocols, such as DNP3, EtherNet/IP, S7comm, and OPC UA, to broaden threat monitoring coverage.
    \item \textbf{Physical Testbed Validation:} Deploying and evaluating the monitoring pipeline within a physical ICS testbed utilizing hardware PLCs and HMIs to validate multi-stage attack chains under realistic operational conditions.
    \item \textbf{Anomaly Detection and Response Automation:} Integrating lightweight machine learning models for anomaly detection and coupling the SIEM layer with Security Orchestration, Automation, and Response (SOAR) platforms to enable safety-bounded, automated threat mitigation.
\end{enumerate}